\section*{异面直线间的距离}

A 组:

\begin{enumerate}

\item 空间和两条异面直线同时都垂直且相交的直线\dotfill{(\kh{B})}

\xx{不一定存在} {有且只有 1 条} {有 1 条或不存在} {有无数条}

\item 在正方体 $ {ABCD} - {A}_{1}{B}_{1}{C}_{1}{D}_{1} $ 中，棱长为 $ a $ ，所在直线与 $ {AD} $ 成异面直线且距离为 $ a $ 的棱有\blkda{4}条.

\item 长方体 $ {{ABCD}} - {{A}}_{1}{{B}}_{1}{{C}}_{1}{{D}}_{1} $ 中， $ {{AB}} $ 和 $ {D{D}_{1}} $ 的公垂线段是\blkda{AD}. $ {A}{{A}}_{1} $ 和 $ {B}{{C}}_{1} $ 的公垂线段是\blkda{AB}.

\item 已知正方体 $ {ABCD} - {A}_{1}{B}_{1}{C}_{1}{D}_{1} $ 的棱长为 $ a $ ，异面直线 $ {DB} $ 与 $ C{C}_{1} $ 之间的距离为\blkda{$\dfrac{\sqrt{2}a}{2}$}.

\item 正方体 $ {{ABCD}} - {{A}}_{1}{{B}}_{1}{{C}}_{1}{{D}}_{1} $ 中，边长为 4，则异面直线 $ {{A}}_{1}{D} $ 与 $ {{B}}_{1}{{D}}_{1} $ 的距离为\blkda{$\dfrac{4\sqrt{3}}{3}$}.

    
\end{enumerate}


B 组:

\begin{enumerate}

\item 空间四边形 $ {{ABCD}} $ 中， $ {{AB}} = {{BD}} = {{AD}} = 2,{{BC}} = {{CD}} = {3} $ ,延长 $ {{BC}} $ 到 $ {E} $ ,使得 $ {{BC}} = {{CE}},{F} $ 为 $ {{BD}} $ 中点,则异面直线 $ {{AF}} $ 和 $ {{DE}} $ 的距离为\blkda{1}.

\item 边长为 1 的两个正方形 $ {ABCD} $ 和 $ {CDEF} $ 构成大小为 $ \dfrac{2}{3}\pi $ 的二面角，则异面直线 $ {AE} $ 和 $ {CD} $ 之间的距离为\blkda{$\dfrac{1}{2}$}.

\item 四面体 $ {SABC} $ 中， $ {SA} = {BC} = {13} $ ， $ {SB} = {AC} = {14} $ ， $ {SC} = {AB} = {15} $ ， 则异面直线 $ {A}{S} $ 与 $ {B}{C} $ 的距离为\blkda{$3\sqrt{14}$}.

\begin{jieda}
    【补形】将四面体 \( {SABC} \) 补成长方体 \( {SDBG} - {EAFC} \) ,连接 \( {DE} \) 交 \( {AS} \) 于点 \( M \) ,连接 \( {FG} \) 交 \( {BC} \) 于点 \( N \) ,连接 \( {MN} \) ,
则 \( M \) 、 \( N \) 分别为 \( {DE} \) 、 \( {BC} \) 的中点,


    \begin{figure}[h]
        \centering
        \includegraphics[width=0.25\linewidth]{bodyxb//src/auhf'ahfafsfafa.png}
    \end{figure}
    由已知可得 \( \left\{  \begin{array}{l} S{A}^{2} = S{D}^{2} + S{E}^{2} = {169} \\  S{B}^{2} = S{D}^{2} + S{G}^{2} = {196} \\  S{C}^{2} = S{E}^{2} + S{G}^{2} = {225} \end{array}\right. \) ,可得 \( {SG} = 3\sqrt{14} \) ,
因为 \( {BD}//{CE} \) 且 \( {BD} = {CE} \) ,故四边形 \( {BDE}C \) 为平行四边形,则 \( {BC}//{DE} \) 且 \( {BC} = {DE} \) ,
又因为 \( M \) 、 \( N \) 分别为 \( {DE} \) 、 \( {BC} \) 的中点,所以, \( {DM}//{BN} \) 且 \( {DM} = {BN} \) ,
故四边形 \( {BDMN} \) 为平行四边形,故 \( {MN}//{BD} \) 且 \( {MN} = {BD} = {SG} = 3\sqrt{14} \) ,
\( \because {BD} \bot \) 平面 \( {SDAE},{AS} \subset \) 平面 \( {SDAE},\therefore {BD} \bot  {AS} \) ,即 \( {MN} \bot  {AS} \) ,
同理可得 \( {MN}\bot {BC} \) ,故异面直线 \( {AS} \) 与 \( {BC} \) 的距离为 \( {MN} = 3\sqrt{14} \) .
故答案为: \( 3\sqrt{14} \) .

\end{jieda}


\item 如图，在棱长为 2 的正方体 $ {ABCD} - {A}_{1}{B}_{1}{C}_{1}{D}_{1} $ 中， $ E $ 为 $ {BC} $ 的中点,点 $ P $ 在线段 $ {D}_{1}E $ 上,点 $ P $ 到直线 $ C{C}_{1} $ 的距离的最小值为\blkda{$\dfrac{2\sqrt{5}}{5}$}.

\begin{center}

    \includegraphics[width=0.25\linewidth]{bodyxb//src/gsddlihjgsoigo'aedgs.png}
\end{center}



\item 已知菱形 $ {ABCD} $ 中， $ \angle {BAD} = \dfrac{\pi }{3} $ ，沿对角线 $ {BD} $ 折起，使二面角 $ A - {BD} - C $ 的平面角为 $ \theta $ ，若异面直线 $ {AC} $ 与 $ {BD} $ 的距离是菱形边长的 $ \dfrac{3}{4} $ ，则 $ \theta  = $ \blkda{$\dfrac{\pi}{3}$}.

\begin{jieda}

如图,设菱形的边长为 \( {2a} \) ,连接两条对角线 \( {AC} \cap  {BD} = O \)
易得 \( {AC} \bot  {BD},{AO} = {OC} = \sqrt{3}a \)

    \begin{figure}[h]
        \centering
        \includegraphics[width=0.25\linewidth]{bodyxb//src/aihfahfoaofhafa.png}
    \end{figure}


    菱形 \( {ABCD} \) 沿对角线 \( {BD} \) 折起,连接 \( {AC} \) ,得到三棱锥 \( A - {BCD} \)
在菱形 \( {ABCD} \) 中, \( {AC} \bot  {BD} \) ,翻着后垂直不变,即 \( \left\{  \begin{array}{l} {AO} \bot  {BD} \\  {CO} \bot  {BD} \end{array}\right. \)
所以 \( \angle {AOC} \) 是二面角 \( A - {BD} - C \) 的平面角,即 \( \angle {AOC} = \theta \)
又因为 \( \left\{  \begin{array}{l} {AO} \bot  {BD} \\  {CO} \bot  {BD} \\  {AO} \cap  {CO} = O \\  {AO},{CO} \subset  \text{ 面 }{AOC} \end{array}\right. \)
所以 \( {BD} \bot \) 平面 \( {AOC} \) ,取 \( {AC} \) 中点 \( M \) ,连接 \( {OM} \)
又因为 \( {OM} \subset \) 平面 \( {AOC} \)
所以 \( {OM} \bot  {BD} \)
在 \( \bigtriangleup  {AOC} \) 中, \( {AO} = {OC} = \sqrt{3}a \) ,并且 \( M \) 为 \( {AC} \) 的中点,
所以 \( {OM} \bot  {AC} \)
故 \( {OM} \) 是异面直线 \( {AC} \) 与 \( {BD} \) 的距离
又因为异面直线 \( {AC} \) 与 \( {BD} \) 的距离是菱形边长的 \( \frac{3}{4} \)
所以 \( {OM} = \frac{3}{4} \times  {2a} = \frac{3}{2}a \)
在 \( \mathrm{R}{t\Delta AOM} \) 中, \( \angle {AOM} = \frac{\theta }{2} \)
所以 \( \cos \frac{\theta }{2} = \frac{OM}{AO} = \frac{\frac{3}{2}a}{\sqrt{3}a} = \frac{\sqrt{3}}{2} \) ,又因为 \( \frac{\theta }{2} \in  \left( {0,\frac{\pi }{2}}\right) \)
所以 \( \frac{\theta }{2} = \frac{\pi }{6}\therefore \theta  = \frac{\pi }{3} \)

\end{jieda}


\item 下列命题其中是真命题的有 \blkda{(2)(3)}

\begin{enumerate}
    \item 两条异面直线的公垂线有无数条;
    \item 异面直线之间的距离就是两条异面直线上点之间距离的最小值;
    \item 过两条异面直线中的一条有且只有一个平面与另一条直线平行.
\end{enumerate}



\item 正方体 $ {{ABCD}} - {{A}}_{1}{{B}}_{1}{{C}}_{1}{{D}}_{1} $ 中，边长为 1，与直线 $ {{A}}_{1}{D} $ 异面且距离为 $ \dfrac{\sqrt{3}}{3} $ 的面对角线有 \blkda{4}条.



\item 若直线 $ a $ 和已知直线 $ b $ 满足：① $ a $ 、 $ b $ 是异面直线，② $ a $ 、 $ b $ 之间的距离是定值，③a与 $ b $ 所成的角也是定值，则直线 $ a $ \dotfill{(\kh{D})}

\xx{仅有一条}{有两条} {有四条}{有无数条}

\item 空间四边形 $ {{ABCD}} $ 中， $ {{AB}} = {{CD}} = {4},{{BC}} = {{AD}} = {3},{{AD}} \bot  {{DC}} $ , $ {{AD}} \bot  {{BC}},{{AB}} \bot  {{BC}} $ ,求异面直线 $ {{AD}} $ 和 $ {{BC}} $ 的距离.


\begin{jieda}
    【鳖臑，补形】$\sqrt{7}$

    \begin{figure}[h]
        \centering
        \includegraphics[width=0.25\linewidth]{bodyxb//src/gsdijgoahdfgohahf.png}
    \end{figure}
\end{jieda}


\item 所有棱长都为 1 的四面体 $ {{ABCD}} $ 中，找到异面直线 $ {{AB}} $ 和 $ {{CD}} $ 的公垂线， 求出 $ {{AB}} $ 和 $ {{CD}} $ 的距离.
\begin{jieda}

如图,取 \( {AB} \) 中点 \( E,{CD} \) 中点 \( F \) ,连接 \( {BF},{AF},{EF} \) ,
由已知 \( {AF} = {BF} = \sqrt{{1}^{2} - {\left( \frac{1}{2}\right) }^{2}} = \frac{\sqrt{3}}{2}，\therefore {EF} \bot  {AB} \) ,同理 \( {EF} \bot  {CD} \) ,
所以 \( {EF} \) 是 \( {AB} \) 与 \( {CD} \) 的公垂线, \( {AB} \) 与 \( {CD} \) 的距离即为线段 \( {EF} \) 的长.
且 \( {EF} = \sqrt{A{F}^{2} - A{E}^{2}} = \sqrt{\frac{3}{4} - \frac{1}{4}} = \frac{\sqrt{2}}{2} \) ,
所以 \( {AB} \) 与 \( {CD} \) 的距离即为 \( \frac{\sqrt{2}}{2} \) .

    \begin{figure}[h]
        \centering
        \includegraphics[width=0.25\linewidth]{bodyxb//src/fa'hfaouhfiahfouha.png}
    \end{figure}
\end{jieda}

\item 正方体 $ {{ABCD}} - {{A}}_{1}{{B}}_{1}{{C}}_{1}{{D}}_{1} $ 中，边长为 2，求异面直线 $ {{CD}} $ 与 $ {A}{{C}}_{1} $ 的距离.

\begin{jieda}
    异面直线的距离即为BC到面$ABC_1D_1$的距离，答案为$\sqrt{2}$。
\end{jieda}

\item 已知 $ {S} $ 是矩形 $ {{ABCD}} $ 所在平面外一点， $ {{SA}} \bot  {{BC}},{{BS}} \bot  {{CD}},{{SA}} $ 与 $ {{CD}} $ 所成角的大小为 $ \dfrac{\pi }{3},{{SD}} $ 与 $ {{BC}} $ 所成角大小为 $ \dfrac{\pi }{6},{{SA}} = {1} $ ,分别求直线 $ {{SA}} $ 与 $ {{CD}} $ 的距离及 $ {{SB}} $ 与 $ {{AD}} $ 的距离.

\begin{jieda}

\( \because {AD}//{BC},{AB}//{CD},{SA}\bot {BC},{BS}\bot {CD}\therefore {SA}\bot {AD},{SB}\bot {AB} \)
因为 \( {SD} \) 与 \( {BC} \) 所成角大小为 \( \frac{\pi }{6} \) ,而 \( {BC}//{AD} \) ,则 \( \angle {SDA} = \frac{\pi }{6} \) ,
因为 \( {SA} \) 与 \( {CD} \) 所成角大小为 \( \frac{\pi }{3} \) ,而 \( {AB}//{CD} \) ,则 \( \angle {SAB} = \frac{\pi }{3} \) ,
\( {SA} = 1 \) ,则 \( {SD} = 2,{AD} = \sqrt{3},{AB} = \frac{1}{2} \) ,
又 \( {ABCD} \) 是矩形,
所以线段 \( {AD} \) 是直线 \( {SA} \) 与 \( {CD} \) 的公垂线段,线段 \( {AB} \) 是 \( {SB} \) 与 \( {AD} \) 的公垂线线段,
所以直线 \( {SA} \) 与 \( {CD} \) 的距离是 \( \sqrt{3},{SB} \) 与 \( {AD} \) 的距离是 \( \frac{1}{2} \) .
    \begin{figure}[h]
        \centering
        \includegraphics[width=0.25\linewidth]{bodyxb//src/sdoufhauhfOAuhwOFa.png}
    \end{figure}
\end{jieda}


\item 四棱锥 $ P - {ABCD} $ 中， $ {PA}\bot $ 平面 $ {ABCD} $ ，底面 $ {ABCD} $ 是平行四边形,且 $ \angle {BAD} = \dfrac{2\pi }{3},{AB} = 2,{PA} = 3,{AD} = 4,M $ 是 $ {BC} $ 的中点.

\begin{center}

    \includegraphics[width=0.25\linewidth]{bodyxb//src/dig'sodihgsohdgagfa.png}
\end{center}

\begin{enumerate}
    \item 求二面角 $ A - {PM} - B $ 的余弦值；
    \item 求异面直线 $ {AD} $ 和 $ {PM} $ 之间的距离.
\end{enumerate}


\begin{jieda}

\begin{figure}[h]
    \centering
    \includegraphics[width=0.25\linewidth]{bodyxb//src/lfhajhfAOHFwfa.png}
\end{figure}
    \begin{enumerate}
        \item 取 \( {BM} \) 中点 \( E \) ,连接 \( {AE},{PE} \) ,作 \( {AH}\bot {PE} \) ,垂足为 \( H \) ,
再过点 \( {A作AG} \bot  {PM} \) ,连接 \( {HG} \) ,
根据题意可知 \( \bigtriangleup  {ABM} \) 为正三角形,
则 \( {AE} \bot  {BM},{AE} = \sqrt{3},{AM} = 2,\therefore {PM} = {PB} = \sqrt{13},{PE} \bot  {BM} \) ,
又 \( {PE} \cap  {AE} = E,{PE}\text{、}{AE} \subset \) 平面 \( {PAE} \) ,则 \( {BM}\bot \) 平面 \( {PAE} \) ,
因为 \( {AH} \subset \) 平面 \( {PAE} \) ,则 \( {BM} \bot  {AH} \) ,
又 \( {BM} \cap  {PE} = H,{BM}\text{、}{PE} \subset \) 平面 \( {PBM} \) ,
所以 \( {AH} \bot \) 平面 \( {PBM} \) ,
因为 \( {PM},{HG} \subset \) 平面 \( {PBM} \) ,所以 \( {AH} \bot  {PM},{AH} \bot  {HG} \) ,
因为 \( {AG} \cap  {AH} = A,{AG}\text{、}{AH} \subset \) 平面 \( {AGH} \) ,
所以 \( {PM} \bot \) 平面 \( {AGH} \) ,
因为 \( {HG} \subset \) 平面 \( {AGH} \) ,所以 \( {HG} \bot  {PM} \) ,
所以 \( \angle {AGH} \) 为二面角 \( A - {PM} - B \) 的平面角,
在Rt \( \bigtriangleup  {APM} \) 中, \( {AG} = \frac{{AM} \cdot  {AP}}{PM} = \frac{6}{\sqrt{13}} \) ,
在Rt \( \bigtriangleup {APE} \) 中, \( {AH} = \frac{{AE} \cdot  {AP}}{PE} = \frac{3\sqrt{3}}{2\sqrt{3}} = \frac{3}{2} \) ,
在Rt \( {\Delta AGH} \) 中, \( \cos \angle {AGH} = \frac{GH}{AG} = \frac{\sqrt{A{G}^{2} - A{H}^{2}}}{AG} = \frac{\sqrt{\frac{36}{13} - \frac{9}{4}}}{\frac{6}{\sqrt{13}}} = \frac{\sqrt{3}}{4} \) ,
所以二面角 \( A - {PM} - B \) 的余弦值为 \( \frac{\sqrt{3}}{4} \) .

\item 根据底面 \( {ABCD} \) 是平行四边形,所以 \( {AD}//{BC} \) ,
因为 \( {BC} \subset \) 平面 \( {PBC},{AD} ⊄ \) 平面 \( {PBC} \) ,
故 \( {AD}// \) 平面 \( {PBC} \) ,
所以线段 \( {AH} \) 的长度即为直线 \( {AD} \) 与平面 \( {PBC} \) 间的距离,
也即异面直线 \( {AD} \) 和 \( {PM} \) 之间的距离.
由上可知 \( {AH} = \frac{3}{2} \) ,所以异面直线 \( {AD} \) 和 \( {PM} \) 之间的距离为 \( \frac{3}{2} \) .
    \end{enumerate}

    
\end{jieda}

\end{enumerate}


C 组:
\begin{enumerate}

\item 已知在正方体 $ {ABCD} - {A}_{1}{B}_{1}{C}_{1}{D}_{1} $ 中， $ E,F,G $ 分别是棱 $ {D}_{1}{C}_{1},{AB},{BC} $ 的中点， $ H $ 是棱 $ {CD} $ 上一点，则下列命题中正确的个数为\blkda{3}
\begin{enumerate}
    \item 异面直线 $ {D}_{1}H $ 与 $ {AB} $ 之间的距离为定值；

    \item 平面 $ {EFG}// $ 平面 $ A{D}_{1}H $ ；

    \item 设平面 $ {A}_{1}{AH} \cap $ 平面 $ {A}_{1}{B}_{1}{C}_{1}{D}_{1} = l $ ，则 $ {AH}//l $ ；

    \item 直线 $ {A}_{1}{C}_{1} $ 与平面 $ {AB}{C}_{1}{D}_{1} $ 所成的角为 $ {30}^{ \circ  } $ .
\end{enumerate}

\begin{jieda}
    \begin{figure}[h]
        \centering
        \includegraphics[width=0.25\linewidth]{bodyxb//src/alhfaouhfaohwfawef.png}
    \end{figure}
    对于①,过 \( H \) 点作 \( {AD} \) 的平行线,与 \( {AB} \) 相交于点 \( M \) ,则 \( {HM} \) 为异面直线 \( {D}_{1}H \) 与 \( {AB} \) 的公垂线段,且长度等于正方体棱长,所以异面直线 \( {D}_{1}H \) 与
\( {AB} \) 之间的距离为定值,①正确；
对于②,平面 \( {ABCD} \cap \) 平面 \( {EFG} = {FG} \) ,平面 \( {ABCD} \cap \) 平面 \( A{D}_{1}H = {AH} \) ,若平面 \( {EFG} \parallel \) 平面 \( A{D}_{1}H \) ,则有 \( {FG}//{AH} \) ,而 \( {FG}//{AH} \) 不一定
成立,②错误；
对于③,平面 \( {ABCD}// \) 平面 \( {A}_{1}{B}_{1}{C}_{1}{D}_{1} \) , \( {AH} \subset \) 平面 \( {ABCD} \) ,所以 \( {AH}// \) 平面 \( {A}_{1}{B}_{1}{C}_{1}{D}_{1} \) ,
\( {AH} \subset \) 平面 \( {A}_{1}{AH} \) ,平面 \( {A}_{1}{AH} \cap \) 平面 \( {A}_{1}{B}_{1}{C}_{1}{D}_{1} = l \) ,则 \( {AH}\parallel l \) ,③正确；
对于④, \( N \) 为 \( A{D}_{1} \) 的中点,连接 \( {A}_{1}N,{C}_{1}N \)
平面 \( {AB}{C}_{1}{D}_{1} \bot \) 平面 \( {AD}{D}_{1}{A}_{1} \) ,平面 \( {AB}{C}_{1}{D}_{1} \cap \) 平面 \( {AD}{D}_{1}{A}_{1} = A{D}_{1},{A}_{1}N \subset \) 平面 \( {AD}{D}_{1}{A}_{1},{A}_{1}N \bot  A{D}_{1} \) ,
所以 \( {A}_{1}N \bot \) 平面 \( {AB}{C}_{1}{D}_{1} \) ,直线 \( {A}_{1}{C}_{1} \) 与平面 \( {AB}{C}_{1}{D}_{1} \) 所成的角为 \( \angle {A}_{1}{C}_{1}N \) ,
在Rt \( \Delta {A}_{1}{C}_{1}N \) 中, \( {A}_{1}N = \frac{1}{2}{A}_{1}{C}_{1},\angle {A}_{1}{C}_{1}N = {30}^{ \circ  } \) ,直线 \( {A}_{1}{C}_{1} \) 与平面 \( {AB}{C}_{1}{D}_{1} \) 所成的角为 \( {30}^{ \circ  } \) ,④正确.

\end{jieda}


\item 已知长方体底面边长分别为 2,1 , 正四面体的四个顶点分别在长方体的顶点及长度为 1 的棱所在的直线上, 求长方体的高.

\begin{jieda}

因为任何一个三棱锥均存在唯一的外接平行六面体,只要分别过三组对棱作三组平行平面,
即能围成所要的平行六面体.

    \begin{figure}[h]
        \centering
        \includegraphics[width=0.25\linewidth]{bodyxb//src/alfhaouhfaefae.png}
    \end{figure}


    反之,任意一个平行六面体也一定内接一个四面体,且长方体对应的四面体一定是对棱相等
的四面体,正方体对应的一定是正四面体.
由于一组对棱 \( \left( {C{B}^{\prime }}\right. \) 与 \( A{D}^{\prime } \) )距离是2,即 \( {AB} = 2 \) ,故另一组对棱 \( \left( {AC}\right. \) 与 \( {B}^{\prime }{D}^{\prime } \) )也为2,
即高 \( B{B}^{\prime } = 2 \) .
 
\end{jieda}



\end{enumerate}